\documentclass[10pt]{article}
\usepackage[T1]{fontenc}
\usepackage{lmodern,amsmath,amssymb,amsthm,microtype}
\usepackage[margin=0.82in]{geometry}
\usepackage[colorlinks=true,urlcolor=blue,citecolor=blue]{hyperref}
\newtheorem{theorem}{Theorem}
\newtheorem{corollary}{Corollary}
\title{Squared Digit Sums Are Equidistributed Modulo Their Base}
\author{Qichao Wang\\{\small Hebei University of Technology}\\[0.3em]
Daoyu Dong\\{\small University of Electronic Science and Technology of China}}
\date{October 26, 2026}
\hypersetup{pdftitle={Squared Digit Sums and OEIS A079495},pdfauthor={Qichao Wang and Daoyu Dong}}
\begin{document}
\maketitle
\begin{abstract}
For every integer base $b\ge2$, each residue of the squared digit sum
modulo $b$ has natural density $1/b$. The increasing sequence of
positions of any one residue therefore has asymptotic slope $b$.
This proves the all-base density and slope conjecture stated in
OEIS A079495, including composite bases. We also give exact finite
digit recurrences for every counting function and every sequence
term, and an exact count on ternary length blocks.
\end{abstract}

\section{The question and the full statement}
For the base-$b$ digits of $n\ge0$, let $s_b(n)$ be the sum of their
squares, with $s_b(0)=0$. Write
\[
 C_{b,r}(N)=\#\{0\le n<N:s_b(n)\equiv r\pmod b\},
 \qquad 0\le r<b.
\]
OEIS A079495 \cite{target} lists the positions of residue zero for
$b=3$, starting with zero. Its comments conjecture density $1/b$
and positional slope $b$ in every base. The related digit-sum
entry A079498 \cite{linear} already contains a proof for ordinary
digit sums; the present argument treats squared digits.

\begin{theorem}\label{thm:density}
For every fixed $b\ge2$ and $0\le r<b$,
\[
 C_{b,r}(N)=\frac Nb+o(N)\qquad(N\longrightarrow\infty).
\]
If $a_{b,r}(j)$ is the $j$th nonnegative integer with residue $r$,
indexed by $j\ge1$, then
\[
 \lim_{j\to\infty}\frac{a_{b,r}(j)}j=b.
\]
\end{theorem}
\begin{proof}
Put $\zeta=\exp(2\pi i/b)$ and
\[
 \lambda_\ell=\frac1b\sum_{d=0}^{b-1}\zeta^{\ell d^2},
 \qquad 1\le\ell<b,\qquad
 \rho=\max_{1\le\ell<b}|\lambda_\ell|.
\]
Every summand has modulus one. The terms for $d=0$ and $d=1$
are $1$ and $\zeta^\ell$, which are distinct. Strictness in the
triangle inequality therefore gives $\rho<1$. This uses no
primality assumption on $b$.

For a complete block of $k$ unrestricted digits, character
orthogonality gives the exact count
\begin{equation}\label{eq:fourier}
 C_{b,r}(b^k)=\frac{b^k}{b}
 \left(1+\sum_{\ell=1}^{b-1}\zeta^{-\ell r}\lambda_\ell^k\right).
\end{equation}
Its discrepancy from $b^k/b$ is at most
$(b-1)b^{k-1}\rho^k$. More generally a block
$[A b^k,(A+1)b^k)$ fixes the leading digits of $A$ and has
$k$ unrestricted trailing digits. Adding $s_b(A)$ merely shifts
the residue, so the same discrepancy bound holds for every such
block. For $k=0$ use the bound $(b-1)/b$ directly.

The base-$b$ expansion of $N$ decomposes $[0,N)$ into at most
$b-1$ such blocks of each size $b^k$, $0\le k\le L$, where
$L=\lfloor\log_b N\rfloor$. Consequently
\begin{equation}\label{eq:error}
 \left|C_{b,r}(N)-\frac Nb\right|
 \le\frac{(b-1)^2}{b}\sum_{k=0}^{L}(b\rho)^k,
\end{equation}
where the $k=0$ summand is one, also when $\rho=0$.
If $b\rho<1$ this is bounded; if $b\rho=1$ it is
$O_b(\log(N+1))$; if $b\rho>1$ it is
$O_b(N^{1+\log_b\rho})$. In the last case the exponent is
strictly smaller than one. All three cases prove density $1/b$.

In particular each residue occurs infinitely often. At its $j$th
position, $C_{b,r}(a_{b,r}(j)+1)=j$, so
$j=(a_{b,r}(j)+1)/b+o(a_{b,r}(j)+1)$. Division and inversion
give the asserted slope.
\end{proof}

\section{Exact finite counts and sequence terms}
Let $V_k(r)=C_{b,r}(b^k)$. Entirely in integer arithmetic,
\begin{equation}\label{eq:dp}
 V_0(r)=\begin{cases}1,&r=0,\\0,&r\ne0,\end{cases}
 \qquad
 V_{k+1}(r)=\sum_{d=0}^{b-1}V_k(r-d^2\bmod b).
\end{equation}
To count below an arbitrary $N$, read its digits from left to right.
At a position having digit $c$ and $k$ positions to its right,
let $u$ be the squared sum of its already fixed prefix. For each
$0\le d<c$, add
\[
 V_k(r-u-d^2\bmod b).
\]
Then fix the actual digit $c$ and update $u\leftarrow u+c^2$.
The disjoint first-differing-digit cases partition $[0,N)$, proving
this exact digit formula. Leading zeros cause no change in $s_b$.

These counts also compute every sequence term: take a length $k$
with $V_k(r)\ge j$ and perform binary search for the least $N$
with $C_{b,r}(N)\ge j$; then $a_{b,r}(j)=N-1$.
Theorem~\ref{thm:density} proves that such a finite $k$ always
exists. This is a finite digit algorithm, not an elementary closed
expression for the $j$th position.

\begin{corollary}
For base three, the number of A079495 terms below $3^k$ is
\[
 V_k(0)=
 \begin{cases}
  3^{k-1},&k\text{ odd},\\[2pt]
  \bigl(3^k+2(-3)^{k/2}\bigr)/3,&k\text{ even}.
 \end{cases}
\]
The even case includes $k=0$, when the value is one.
\end{corollary}
\begin{proof}
For $\zeta=\exp(2\pi i/3)$, the digit squares are $0,1,1$
modulo three, so $\lambda_1=(1+2\zeta)/3=i/\sqrt3$ and
$\lambda_2=-i/\sqrt3$. Substitute these values into
\eqref{eq:fourier} with $r=0$.
\end{proof}

\section{Scope and independent checks}
The accompanying \texttt{verify.py} compares every integer-prefix
count for bases 2 through 16 and $0\le N\le5000$ with direct
digit extraction. It compares all 63 displayed A079495 terms,
checks the exact ternary block counts through length 100 by
integer convolution, and independently checks sequence selection
for each residue in each tested base. The all-base and asymptotic
claims follow from the proof, not finite numerical ratios.
No literature priority or OEIS acceptance is asserted.

\begin{thebibliography}{9}
\bibitem{target} C. Alves, OEIS A079495,
\url{https://oeis.org/A079495}, accessed October 26, 2026.
\bibitem{linear} C. Alves, OEIS A079498, including M. Dekking's
April 18, 2019 proof for ordinary digit sums,
\url{https://oeis.org/A079498}, accessed October 26, 2026.
\end{thebibliography}
\end{document}
